
\begin{frame}{Pourquoi pas une unique table\,?}

\begin{small}
\begin{tabular}{c|c|c|c|c}
\textbf{nom} & \textbf{capacité} & \textbf{type} & \textbf{lieu} & \textbf{codeActivité}\\ \hline 
Causses & 45 & Auberge & Cévennes & Randonnée\\ \hline 
Génépi & 134 & Hôtel & Alpes & Piscine\\ \hline 
Génépi & 134 & Hôtel & Alpes & Ski\\ \hline 
U Pinzutu & 10 & Gîte & Corse & Plongée\\ \hline 
U Pinzutu & 10 & Gîte & Corse & Voile\\ \hline 
Tabriz & 34 & Hôtel & Bretagne & Voile\\ \hline 
\end{tabular}
\end{small}

\vfill

Risques d'anomalies et d'incohérences en \alert{insertion}, \alert{modification}
et \alert{suppression}. 
\end{frame}


\begin{frame}{Anomalies d'insertions}

\alert{Pas d'insertion de logement sans connaître au moins une activité}
\vfill

Ou il faut accepter beaucoup de valeurs nulles. Très difficilement gérable.

\vfill

\alert{À chaque nouvelle activité, je dois répéter les informations du logement}
\vfill

Source de redondances et d'incohérences.

\end{frame}


\begin{frame}{Anomalies de mises à jour}

\alert{Si je modifie un logement (capacité) je dois le faire autant de fois qu'il y a d'activités}
\vfill

Si j'oublie, la base devient incohérente

\vfill

\alert{Si je supprime une activité, je risque de supprimer le logement}
\vfill

Il faut donc tester si l'activité supprimée est la dernière du logement

\end{frame}


%%%%%%%%%%%%%%%
\begin{frame}{Schémas normalisés}

Les anomalies précédentes seraient  difficilement acceptables.

\vfill

En pratique, on applique 
un processus de \alert{normalisation}

\begin{redbullet}
  \item L'ensemble des attributs est \alert{décomposé} en plusieurs tables
  \item Cette répartition se fait \alert{sans perte d'information}. 
     Les \alert{jointures} permettent de reconstituer les liens avant décomposition.
     
   \item L'information initiale peut alors être présentée sous  forme de \alert{vue},
     sans les inconvénients de la matérialisation.
\end{redbullet}
 
 \vfill
 La conception d'une base relationnelle consiste à obtenir un schéma normalisé
 et sans perte d'information. 
\end{frame}



%%%%%%ù

\begin{frame}{À retenir}

On s'applique à obtenir des schémas normalisés, sans anomalie
\vfill
\begin{redbullet}
  \item C'est un choix qui simplifie la tâche des applications
  \item En contrepartie il implique de reconstituer, par jointure, l'information
     quand elle est répartie dans plusieurs tables.
   \item L'information de départ peut être représentée par une vue
       (\alert{calculée}), sans les anomalies d'une table (\alert{matérialisée})
 \end{redbullet}

\vfill

Se reporter au premier cours sur le modèle relationnel (notions
de vues matérialisées ou non)
\end{frame}

